\section{Better, Not Just Less Bad}
\label{sec:12.1}
What would it mean for AI to make a life better, not just avoid making it worse? The economist Amartya Sen's capability approach gives the question a definite shape: well-being is not the resources a person has but the real freedom to do and be what they have reason to value \autocite{sen1999development}. Two systems can hand someone the same resource and deliver different freedom. A personalized tutor that adapts to how a specific student learns expands a capability; the same technology, gating what a student is allowed to study against a government's list of approved subjects, does not, no matter how well it teaches within those limits. Which of the two a given deployment turns out to be is decided by what makes a system resistant to capture rather than merely capable.

The approach arrives with a commitment the tutor example does not show, and section~\ref{sec:2.3.1} takes up the half of it that bears on who counts. What bears here is narrower. Martha Nussbaum built her list of central capabilities partly against the case of disability, naming it the first of three frontiers where a theory grounding justice in mutual advantage among roughly equal parties cannot say what is owed \autocite{nussbaum2006frontiers}, and the systems in question sort people by inference from behavior, where atypical behavior is what they are worst at.

There is a worked case. A video-interview product scored candidates on facial expression for years, against sustained objection that the technique was never a validated predictor and that scoring a face and a voice screens out candidates with disabilities such as autism or speech impairments \autocite{maurer2021hirevue,doj2022algorithms}. The scoring was dropped on an unrelated ground: the vendor found the nonverbal signal contributed about 0.25 percent to its models' predictive power \autocite{kahn2021hirevue}. An independent algorithmic audit published ten months later cannot have caused that, and says so. What it examined was one representative pre-built assessment for early-career candidates, and it contains no analysis of the system's training data or code, so its conclusion that the assessments “work as advertised with regard to fairness and bias issues” is a finding about that assessment \autocite{johnson2021auditing}.

Read as an accuracy problem, that is a classifier wanting better training data. Read as a capability question, it is a system narrowing the range of lives available to a group of people by measuring them against a norm they were never going to match — and neither a more accurate classifier nor an audit scoped that way would have reached it.

Taiwan's vTaiwan process is the clearest existing case of AI-assisted tooling expanding a capability, public deliberation, rather than substituting for it. A 2026 project led by the University of Chicago with Princeton and SRI International runs the same logic against censorship, using AI to distinguish deliberate state interference with online information from ordinary network failure at population scale \autocite{uchicago2026censorship}. Neither tool resists authoritarianism by itself. Both are evidence that the same capacity — modeling patterns in enormous, noisy data — can be pointed at expanding what a person can do rather than at narrowing it.

That symmetry is the whole argument for antifascist design rather than a general presumption that AI helps. A capability expanded is not guaranteed; it has to be built in, against a live alternative in which the same underlying capacity gets pointed at surveillance, censorship, or narrowing what a system's users are permitted to want.
